\subsection{问题三的补充假设}
\begin{enumerate}[label=\textbf{假设\arabic*：},leftmargin=4.7em,labelsep=0.4em,itemsep=0.2em]
\item 每日全时段电价在 $0{:}00$ 可获得，购电承诺的调整只影响尚未执行时段；同一时段的最终结算只相对该日 $0{:}00$ 原计划进行一次。
\item 光伏逐小时预报在各发布时点可用，10 分钟预测由当时已观测光伏和后续预报节点因果插值得到；未来实际光伏不参与预测构造。
\item 未来负荷采用此前已结束的最近至多四个同星期日逐时段均值预测；当前执行时段才读取真实负荷，以保持计划阶段的非前瞻性。
\item 储能仍在公共母线侧计量，充、放电效率对称分解为 $\sqrt{0.9}$，不计售电、自放电、辅助功耗和退化费用；外网接口容量未给定，故不额外设购电功率上限。
\item 日末储电量不强制返回日初状态，而以次日因果预测构造的终端价值表示其机会成本；该虚拟价值不计入当天实际购电账单。
\end{enumerate}
